\newgeometry{top=1.5cm,bottom=1.5cm,left=2cm,right=2cm}
\pagestyle{empty}
\AddToShipoutPictureFG{%
  \AtPageUpperLeft{%
    \put(\LenToUnit{19cm},\LenToUnit{-2cm}){\makebox(0,0){\rotatebox{90}{\small\thepage}}}%
  }%
}
\begin{landscape}
\makeatletter
\chaptermark{MATRIZ DE CONSISTENCIA}
\addcontentsline{toc}{chapter}{MATRIZ DE CONSISTENCIA}
{\centering\bfseries\normalsize MATRIZ DE CONSISTENCIA\par}
\vspace{6pt}
\begin{spacing}{1}
\footnotesize
\setcounter{chapter}{1}
\setcounter{table}{0}
\setlength{\tabcolsep}{4pt}
\begin{longtable}{@{}p{4.9cm}p{4.9cm}p{4.9cm}p{4.9cm}p{4.9cm}@{}}
\multicolumn{3}{@{}l@{}}{\textit{\scriptsize UNIVERSIDAD NACIONAL DE INGENIERÍA}} &
\multicolumn{2}{@{}r@{}}{\textit{\scriptsize MATRIZ DE CONSISTENCIA}} \\[-2pt]
\multicolumn{3}{@{}l@{}}{\textit{\scriptsize \@facultad}} &
\multicolumn{2}{@{}r@{}}{} \\
\multicolumn{5}{@{}p{24.5cm}@{}}{\hrulefill}\\[2pt]
\caption{Matriz de Consistencia}\label{tab:matriz}\\
\toprule
Problema & Objetivos & Hipótesis & Variables e Indicadores & Metodología \\
\midrule
\endfirsthead
\multicolumn{3}{@{}l@{}}{\textit{\scriptsize UNIVERSIDAD NACIONAL DE INGENIERÍA}} &
\multicolumn{2}{@{}r@{}}{\textit{\scriptsize MATRIZ DE CONSISTENCIA}} \\[-2pt]
\multicolumn{3}{@{}l@{}}{\textit{\scriptsize \@facultad}} &
\multicolumn{2}{@{}r@{}}{} \\
\multicolumn{5}{@{}p{24.5cm}@{}}{\hrulefill}\\[2pt]
\multicolumn{5}{l}{\small\itshape (Tabla \ref{tab:matriz} -- continuación)}\\
\toprule
Problema & Objetivos & Hipótesis & Variables e Indicadores & Metodología \\
\midrule
\endhead
\midrule
\multicolumn{5}{r}{\small\itshape Continúa en la página siguiente}\\
\bottomrule
\multicolumn{5}{@{}p{24.5cm}@{}}{\hrulefill}\\[2pt]
\multicolumn{3}{@{}l@{}}{\textit{\scriptsize \@title}} &
\multicolumn{2}{@{}r@{}}{} \\[-2pt]
\multicolumn{3}{@{}l@{}}{\textit{\scriptsize Bach.\ \@author}} &
\multicolumn{2}{@{}r@{}}{} \\
\endfoot
\bottomrule
\multicolumn{5}{@{}p{24.5cm}@{}}{\hrulefill}\\[2pt]
\multicolumn{3}{@{}l@{}}{\textit{\scriptsize \@title}} &
\multicolumn{2}{@{}r@{}}{} \\[-2pt]
\multicolumn{3}{@{}l@{}}{\textit{\scriptsize Bach.\ \@author}} &
\multicolumn{2}{@{}r@{}}{} \\
\endlastfoot
\textbf{Problema General:} \newline
¿En qué medida el uso intensivo e irrestricto de Inteligencia Artificial
Generativa (IAG) en la codificación de aplicaciones empresariales incrementa
la Deuda Técnica Acumulada y degrada la mantenibilidad en el sector de Seguros
y Reaseguros?
&
\textbf{Objetivo General:} \newline
Determinar y cuantificar el impacto del uso de herramientas de IAG en el
incremento de la Deuda Técnica Acumulada y la degradación de la
mantenibilidad en aplicaciones empresariales de Seguros y Reaseguros,
mediante un modelo de evaluación estático y dinámico.
&
\textbf{Hipótesis General ($H_1$):} \newline
El uso intensivo de IAG en el desarrollo de aplicaciones empresariales de
Seguros y Reaseguros incrementa de forma estadísticamente significativa
($\alpha = 0.05$) la Deuda Técnica Acumulada ($TDR \ge 35\%$) y reduce la
mantenibilidad ($MI \le 20$), en comparación con el desarrollo tradicional
sin IAG. \newline
\textbf{Hipótesis Nula ($H_0$):} \newline
No existe una diferencia estadísticamente significativa en la Deuda Técnica
ni en la mantenibilidad entre ambos enfoques.
&
\textbf{Variable Independiente (X):} Nivel de Adopción de IAG en la
Codificación. \newline
Indicadores: \% de líneas aceptadas, tasa de sugerencias. \newline
\textbf{Variable Dependiente 1 ($Y_1$):} Deuda Técnica Acumulada (estática). \newline
Indicadores: Technical Debt Ratio (TDR), Complejidad Ciclomática $v(G)$,
\% de duplicación, CBO. \newline
\textbf{Variable Dependiente 2 ($Y_2$):} Mantenibilidad dinámica. \newline
Indicadores: Maintainability Index (MI), Time to Resolve Change Request
(TTR-CR), Time to Fix Bug (TTFB).
&
\textbf{Tipo y Nivel:} Aplicada, cuantitativa, cuasi-experimental y
explicativa. \newline
\textbf{Diseño:} Cuasi-experimental con grupo control (A -- tradicional) y
grupo experimental (B -- con IAG). \newline
\textbf{Población y Muestra:} Módulos core de Tarificación Actuarial,
Emisión y Cesiones de Reaseguros. \newline
\textbf{Herramientas:} SonarQube, ArchUnit, ESLint, JaCoCo, Jira/GitLab API.
\\
\textbf{Problemas Específicos:} \newline
1. ¿Cómo afecta la IAG a la complejidad ciclomática y cognitiva en módulos
de tarificación y suscripción? \newline
2. ¿En qué grado incrementa la duplicación de código y el acoplamiento
estructural en la integración de Reaseguros? \newline
3. ¿De qué manera la sobre-dependencia de IAG impacta en el tiempo requerido
para resolver cambios y bugs?
&
\textbf{Objetivos Específicos:} \newline
1. Evaluar y comparar la Complejidad Ciclomática y Cognitiva en tarificación
y suscripción. \newline
2. Medir el impacto de la IAG en la duplicación y acoplamiento (CBO) en
componentes de Reaseguro. \newline
3. Cuantificar el esfuerzo temporal (TTR-CR y TTFB) para refactorizar y
corregir bugs.
&
\textbf{Hipótesis Específicas:} \newline
$H_{1.1}$: La Complejidad Ciclomática y Cognitiva en tarificación supera
$v(G) > 10$ en más del 40\% de funciones. \newline
$H_{1.2}$: Los componentes de Reaseguro presentan un incremento de al menos
25\% en duplicación y $CBO \ge 15$. \newline
$H_{1.3}$: El tiempo de resolución (TTR-CR y TTFB) es un 30\% mayor debido a
la baja cohesión.
&
\textbf{Variables Intervinientes / de Control:} \newline
-- Experiencia del desarrollador (años en el sector asegurador/IT). \newline
-- Complejidad de la regla de negocio (Suscripción vs. Tarificación vs.
Reaseguro).
&
\textbf{Procesamiento Estadístico:} \newline
-- Pruebas de normalidad (Shapiro-Wilk). \newline
-- Contrastación con t-Student para muestras independientes o U de
Mann-Whitney ($\alpha = 0.05$).
\\
\end{longtable}
\setlength{\tabcolsep}{6pt}
\normalsize
\end{spacing}
\makeatother
\end{landscape}
\ClearShipoutPictureFG
\pagestyle{fancy}
\restoregeometry
